% Part 1, Problem 4: sqrt{a} irrational for non-square a - Proof by Contradiction
% Easy

\begin{problem}[Easy]
Let $a$ be a positive integer. If $a$ is not a perfect square, prove that $\sqrt{a}$ is irrational.
\end{problem}

\begin{solution}
\textbf{Proof by Contradiction}

\textbf{Step 1: Assume the negation}

Assume, for contradiction, that $\sqrt{a}$ is rational.
Then we can write
\[
\sqrt{a} = \frac{p}{q}
\]
where $p, q \in \mathbb{Z}$, $q \neq 0$, and $\gcd(p,q) = 1$.

\textbf{Step 2: Square both sides}

\begin{align*}
a &= \frac{p^2}{q^2} \\
p^2 &= aq^2
\end{align*}

\textbf{Step 3: Compare prime factorizations}

By the Fundamental Theorem of Arithmetic, every positive integer has a unique prime factorization.

When an integer is squared, every prime exponent in its factorization becomes even.
So:
\begin{itemize}
    \item in $p^2$, every prime appears with an even exponent;
    \item in $q^2$, every prime appears with an even exponent.
\end{itemize}

Now the equation
\[
p^2 = aq^2
\]
shows that $aq^2$ must also have only even prime exponents.

Since $q^2$ already contributes only even exponents, this is possible only if every prime appearing in $a$ also has an even exponent.

\textbf{Step 4: Derive the contradiction}

But if every prime in the factorization of $a$ has an even exponent, then $a$ is a perfect square.

This contradicts the hypothesis that $a$ is \emph{not} a perfect square.

\textbf{Conclusion}

Therefore, our assumption was false, and $\sqrt{a}$ must be irrational.

\hfill $\square$
\end{solution}

\begin{takeaways}
\begin{itemize}
\item \textbf{Parity of Prime Exponents:} A perfect square has only even exponents in its prime factorization
\item \textbf{Why the Method Works:} Squaring forces all prime exponents in $p^2$ and $q^2$ to be even, so the same must be true for $a$
\item \textbf{Generalization:} This extends proofs like ``$\sqrt{2}$ is irrational'' or ``$\sqrt{23}$ is irrational'' to any positive integer that is not a square
\item \textbf{Key Tool:} The argument relies on the Fundamental Theorem of Arithmetic and uniqueness of prime factorization
\end{itemize}
\end{takeaways}
